\scanpage{135}
\mspage{84}

$T_1$ est une sous-catégorie pleine de $T$, $g$ la flèche
d'inclusion. Notons à ce sujet la proposition
suivante

\uline{Proposition} 7 Conditions équivalentes sur l'ouvert
$U$ d'un objet $X$ de $\mathrm{Cat}$ (moyennant $\mathrm W(7\ \mathrm{bis})$) :

a) l'inclusion $i_0:U\to\overline U$ est \deleted{dans $\mathcal W$}\inserted{$\mathcal W$-asphérique},
\deleted{\unsure{b)}} i.e. pour tout $t\in\dot U=\overline U\setminus U$, la catégorie
$U_{/t}$ est $\mathcal W$-asphérique.

b) L'inclusion $i_0:U\to\overline U$ est dans $\mathcal W$, et cette
condition reste satisfaite, quand on remplace
$(X,U)$ par $(X_1,U)$, où $(X_1,U)$ est déduit de
$(X,U)$ au moyen d'une flèche $T_1\to T$ $(=X\setminus U)$
dans $\mathrm{Cat}$, comme rappelé ci-dessus.

\deleted{b') \unsure{La condt} b), \unsure{est remplie|où on remplace} dans le cas où
$X_1$ est une sous-catégorie pleine de $X$ contenant $U$.}

b') Comme b), dans le cas où $T_1$ est la cat. ponctuelle $e$, et
\inserted{où de plus son image est dans $\dot U$.}

\deleted{b'') Comme b'), dans le cas où la sous-catégorie
pleine est de la forme $U\cup\{t\}$, avec $t\in\dot U$.}

\uline{Démonstration} : Il est évident que si dans b, b', b''
on remplace la condition ``$i_0\in\mathcal W$'' par \inserted{la condition plus forte} ``$i_0$ est
$\mathcal W$-asphérique'', les quatre conditions \deleted{\unsure{a) à}} b'') sont équivalentes. Il reste à prouver que
\unsure{la|les} condition\unsure{s} que les $U_{/t}$ soient $\mathcal W$-asphériques
($t\in\overline U$ -- mais c'est trivial pour $t\in U$, donc il
suffit de l'exprimer pour $t\in\overline U\setminus U=\dot U$)
est déjà impliquée par \deleted{\unsure{chacune des}}\inserted{la} plus faible, soit \unsure{b'} \deleted{\unsure{b''}},
et on devine qu'on a ceci :

\deleted{\uline{Corollaire 1} Soit $X_1\unsure{(t)|(b)}$ la sous-catégorie pleine}

% archive sheet number: 135
